\usection{Актуальность исследования}

Подавляющее большинство используемых сегодня языков
программирования не подходит для написания корректного по
построению кода. В случае же, когда требуется
удостовериться в корректном исполнении программы
(исполнение без ошибок, соответствие спецификации и так
далее), программисту необходимо настраивать дополнительные
инструменты: статические анализаторы, автоматическое
тестирование, средства проверки модели и так далее. К
сожалению, ни одно из используемых средств не решает
задачу полностью: какие-то из них не гарантируют
корректности\footnote{Здесь уместно процитировать Эдсгера
В. Дейкстру: ``... тестирование программы может весьма
эффективно продемонстрировать наличие багов, но совершенно
неприменимо, чтобы показать их отсутствие''
\cite{dijkstra1972humble}.}, другие имеют
(сверх)экспоненциальные асимптотики исполнения и не могут
быть использованы для полного анализа проверяемых
систем\footnote{Однако, например, проверку моделей вполне
успешно используют для верификации критических компонентов
системы и рискованных оптимизаций. \cite{amazon2015}}.
В конце концов, сама необходимость дополнительной
настройки, коммерческое лицензирование и сложность
использования мешают распространению подобных технологий.

Так, исчерпывающим решением было бы легковесное,
бесплатное, удобное для использования, встроенное в
стандартную поставку языка программирования средство
формальной спецификации и формальной верификации
программной системы; по определению, сформулированному
Бенджамином К. Пирсом \cite{пирс2012типы}, таким средством
является \textit{система типов}. Приведём его здесь:

\textit{Система типов} --- это гибко управляемый
синтаксический метод доказательства отсутствия в программе
определенных видов поведения при помощи классификации
выражений языка по разновидностям вычисляемых ими
значений (\textit{типам}).

Становится возможным поставить следующую
\textit{исследовательскую проблему}: \\
Какой должна быть система типов языка программирования,
чтобы
\begin{itemize}
\item Она позволяла легко формулировать и легко (в идеале
--- автоматически) доказывать свойства программ,
интересные в мейнстримном программировании?
\item Полученным языком было удобно пользоваться широкому
кругу пользователей?
\end{itemize}

\usection{Степень разработанности исследовательской
проблемы}

К подтверждению гипотезы, что теорию типов можно
использовать для формальной верификации компьютерных
программ, логики и исследователи компьютерных наук
приходили постепенно. Пожалуй, первым удачным
экспериментом в этой области следует считать Исчисление
Конструкций (Calculus of Constructions, CoC)
\cite{COQUAND198895} Тьерри Коканда, в котором можно не
только реализовывать алгоритмы в виде программ на
некотором функциональном языке, но и доказывать их
свойства в той же самой системе; практическая реализация
этого исчисления сегодня известна как Rocq (ранее Coq)
\cite{the_coq_development_team_2024_14542673}.

Эта замечательная особенность CoC возникла в результате
принятия парадигмы \textit{утверждения в качестве
(некоторых) типов} \cite{wadler2015}. В её рамках
доказываемое утверждение $\varphi$ кодируется в виде
некоторого \textit{типа} $F$, представимого в выбранной
системе (так что язык типов должен быть достаточно
выразительным, чтобы кодировать все утверждения, которые
мы хотим доказать). Тогда задача доказательства
$\vdash\varphi$ сводится к задаче построения выражения
$f$, имеющего тип $F$ ($\vdash f:F$); поиск такого
выражения может производиться как вручную, так и с помощью
автоматических средств\footnote{В последнее время стала
популярной методика, по которой доказательство
математических утверждений ищут с помощью больших языковых
моделей, генерирующих программы в одной из таких систем.
\cite{song2025leancopilotlargelanguage}}.

Однако следует иметь в виду, что Rocq и другие системы
интерактивных доказательств в первую очередь являются
\textit{системами доказательств}, нежели \textit{языками
программирования}: как правило, скрипты для таких систем
не запускаются в качестве самостоятельных программ, а
только подаются на вход алгоритму проверки
типов\footnote{В Rocq, однако, есть обходной механизм
\textit{извлечения доказательства}, транслирующий скрипт
Rocq в программу на OCaml/Haskell/Scheme/...}. Значимыми
исключениями являются языки Agda и Lean: в соответствующих
компиляторах есть опции, позволяющие компилировать и
запускать скрипты на этих языках.

Agda \cite{norell:thesis}, как и язык, на котором она
реализована (Haskell), в первую очередь является тестовой
площадкой для различных расширений CoC: доступны опции для
подключения коиндуктивных типов \cite{AltenkirchD10} и
ограниченной рекурсии
\cite{clouston2015programmingreasoningguardedrecursion},
кубических типов \cite{vezzosi2019cubical}, логически
противоречивого правила Type-in-Type и многих других.
К сожалению, это делает ядро системы перегруженным и
малоприменимым в критических ситуациях.

Система Lean \cite{moura2021lean} получила широкую
известность ввиду прорывного по тогдашним меркам
инструментария и распространения в среде
математиков-энтузиастов. Кроме этого, архитектура системы
типов и стандартных библиотек была спроектирована с учётом
удобства переноса доказательств из классической
математики; в решении, избранном создателями Lean, для
этого пришлось пожертвовать достаточно важными свойствами
базовой системы CoC, в частности свойством
\textit{сохранения типизации}: если $e_1:T$ и $e_1\to e_2$
(<<$e_1$ за один шаг вычисляется в $e_2$>>), то $e_2:T$.
Повторюсь: это свойство в Lean НЕ выполняется. Если быть
точным, можно подобрать типизируемый терм $e_1$ такой, что
результат вычисления $e_2$ не будет типизироваться. Данное
обстоятельство, конечно, не делает данную логическую
систему противоречивой (в частности, существует
теоретико-модельное доказательство непротиворечивости
системы Lean \cite{digama2019leantypetheory}), но является
серьёзным неудобством при неумелой работе с системой.
Более того, если такое базовое синтаксическое свойство не
выполнено, возникают сомнения в применимости подобной
системы для программирования, т.е. акта дизайна программ
для исполнения на переписывающей системе.

Таким образом, можно сделать вывод, что на данный момент
не существует \textit{языка программирования} на основе
CoC, сохраняющего важные вычислительные свойства и
пригодного для разработки производительных программных
систем целиком в этом языке.

\usection{Объект и предмет исследования}

\textit{Объектом} данного исследования является система
типов CoC и её расширения, а также практическая реализация
эффективной среды исполнения CoC.

\textit{Предметом} исследования являются такие аспекты CoC
и её расширений, которые сделают её пригодной для
промышленного применения в качестве системы исполнения
высокопроизводительных программ, корректных по построению.

\usection{Цель и задачи исследования}

\textit{Целью} данной работы является поиск системы типов,
основанной на CoC, пригодной для написания корректных по
построению программ и доступной широкому кругу
специалистов в сфере информационных технологий, а также её
эффективная реализация.

Для достижения этой цели необходимо выполнить следующие
\textit{задачи}:

\begin{enumerate}
    \item Исследование возможных (как существующих, так и
        новых) расширений системы CoC;
    \item Отбор наиболее необходимых для мейнстримного
        программирования расширений;
    \item Доказательство сохранения важных вычислительных
        свойств в расширенной системе;
    \item Дизайн эргономичного алгоритма вывода типов,
        поддерживающего существующие и будущие паттерны
        программирования с зависимыми типами;
    \item Реализация эффективной среды исполнения и
        логического ядра системы;
    \item Создание демонстрационных проектов на новом
        языке, показывающих состоятельность выбранного
        подхода.
\end{enumerate}

\usection{Методология и методы исследования}

В силу природы поставленных задач, настоящее исследование
производится на стыке математической логики (теория типов,
теория моделей), алгебры (символьные вычисления,
алгебраическая оптимизация) и теоретической информатики
(теория алгоритмов, теория сложности) с акцентом на
практической реализации. Для каждой из поставленных задач
используются свои методы исследования, в частности:

\begin{enumerate}
    \item Полученная система типов описывается в
        декларативной форме, в виде генценовского
        исчисления секвенций специального вида;
    \item Метатеоретические свойства системы доказываются
        относительно сформулированного исчисления
        структурной индукцией по дереву естественного
        вывода в системе;
    \item Алгоритм вывода типов также формулируется с
        помощью исчисления секвенций, но уже в
        алгоритмической форме, демонстрируется
        корректность вывода относительно базовой системы;
    \item Реализация компилятора и среды исполнения
        опирается на новейшие практики реализации подобных
        программных систем;
    \item Эффективность среды исполнения проверяется
        посредством бенчмаркинга и сравнения
        производительности по сравнению с аналогичными
        средами на эталонных программах, реализованных на
        различных языках.
\end{enumerate}

\usection{Научная новизна исследования}

Данное исследование предлагает новые дизайны для систем
типов с зависимыми типами с акцентом на практическую
применимость в промышленном программировании, основываясь
на новейших результатах в области оснований математики;
кроме этого, настоящая работа содержит значимые
продвижения в техниках создания эффективных сред
исполнения для языков с зависимыми типами, а также в
техниках использования зависимых типов в промышленном
программировании.

\usection{Практическая значимость работы}

Существование удобного в использовании языка
программирования с эффективной средой исполнения,
позволяющего создавать корректные по построению программы,
качественно изменит подход к разработке программного
обеспечения. На подобные программные решения существует
большой спрос, что показывает пример Rust
\cite{nsa2023rust}; технология, позволяющая раздвинуть
границы возможного ещё дальше, будет востребована не
меньше.

\usection{Степень достоверности и апробация результатов}

В лучших традициях дисциплины \textit{теории типов},
достоверность ключевых теоретических результатов,
полученных в ходе исследования, гарантируется
формализацией доказательств в среде интерактивных
доказательств Rocq. Апробация практических результатов
произведена с помощью реализации ряда демонстрационных
проектов в полученном языке программирования,
показывающих состоятельность и перспективность выбранного
подхода.
